% Authority: exact-scope leaf 18; full PDF page 97; printed p.80.
% U:p080.r28.eq.018b | continuation; R28 eq.018b visually replayed
\par\noindent
$h'$, $h''$, \ldots{} sont en nombre pair ou impair. On peut énoncer ce résultat d'une manière un
peu différente, en disant que $\leg{t}{p}$ est égal au produit des seconds membres des relations
$(\varepsilon)$.
% U:p080.r28.par.023 | paragraph-start; R28 par.023 visually replayed
\par\indent
On peut, dans les relations $(\varepsilon)$, changer partout $-a$ en $a$, pourvu qu'en même temps
on change le signe des seconds membres de celles de ces relations où il se trouve un nombre
premier $g$, $g'$, $g''$, \ldots{}, $h$, $h'$, $h''$, \ldots{} de la forme $4n+3$. On aura ainsi:
% U:p080.r28.eq.019 | continuation; R28 eq.019 visually replayed
\begin{equation*}\tag{$\zeta$}
\left\{
\begin{array}{llll}
\leg{a}{g}=\pm1,& \leg{a}{g'}=\pm1,& \leg{a}{g''}=\pm1,& \text{etc.}\\[0.3em]
\leg{a}{h}=\pm1,& \leg{a}{h'}=\pm1,& \leg{a}{h''}=\pm1,& \text{etc.}
\end{array}
\right.
\end{equation*}
% U:p080.r28.par.024 | paragraph-start; R28 par.024 visually replayed
\par\indent
En comparant le second membre de chacune de ces relations avec le second membre de la relation
correspondante du tableau $(\varepsilon)$, on trouvera autant de changements de signe qu'il y a de
nombres premiers $4n+3$ parmi les nombres $g$, $g'$, $g''$, \ldots{}, $h$, $h'$, $h''$, \ldots{}.
% U:p080.r28.par.025 | paragraph-start; R28 par.025 visually replayed
\par\indent
Le nombre des changements sera donc pair, lorsque parmi les nombres $g$, $g'$, $g''$, \ldots{},
$h$, $h'$, $h''$, \ldots{} il s'en trouve un nombre pair de la forme $4n+3$. Dans ce cas qui aura
lieu toutes les fois que $t=gg'g''\ldots{}\times hh'h''\ldots$ est de la forme $4n+1$, le produit
des seconds membres des relations $(\zeta)$ sera donc le même que le produit des seconds membres
des relations $(\varepsilon)$.
% U:p080.r28.par.026 | paragraph-start; R28 par.026 visually replayed
\par\indent
D'un autre côté, on peut, en vertu de la loi de réciprocité, renverser les premiers membres des
relations $(\zeta)$, sans qu'il soit nécessaire de changer le signe d'aucun des seconds membres de
ces relations. Le produit de toutes les expressions renversées, c'est-à-dire
$\leg{gg'g''\ldots{}\times hh'h''\ldots}{a}$ ou $\leg{t}{a}$, est donc égal au produit des seconds
membres des relations $(\zeta)$ et par conséquent aussi, d'après ce qu'on a vu précédemment, égal
au produit des seconds membres des relations $(\varepsilon)$. Or, ayant prouvé plus haut que ce
dernier produit est égal à $\leg{t}{p}$, on a:
% U:p080.r28.eq.020 | continuation; R28 eq.020 visually replayed
\[
\leg{t}{p}=\leg{t}{a}.
\]
% U:p080.r28.par.027 | paragraph-start; R28 par.027 visually replayed
\par\indent
Ce résultat est relatif au cas où $t$ est de la forme $4n+1$; si l'on examine d'une manière
semblable le cas de $t=4n+3$, on trouvera qu'on a alors:
% U:p080.r28.eq.021 | continuation; R28 eq.021 visually replayed
\[
\leg{t}{p}=-\leg{t}{a}.
\]
